% Source: Weinberg GR, physical PDF pp. 599--606
%         (printed pp. 571--578; boundary pages are partial).
% Coverage: Section 15.9, Newtonian Theory of Small Fluctuations; equations
%           (15.9.1)--(15.9.46), including all unnumbered displays.
% Figures/tables/footnotes: reference notes 152a and 153--155; no figures or
%                          tables.
% Status: source-reviewed and compile-clean.
% Uncertainties: the source's inequalities following (15.9.43) select
%                pressure-dominated modes as growing, opposite to (15.9.23)
%                and to the later criterion after (15.9.46); SOURCE ERRATUM
%                corrections are used below.
% MODERNIZATION: R_{\mathrm{source}}(t) -> a(t), with H=\dot a/a. To avoid
% collision with the spatial-curvature constant k, the source's physical
% wavevector k=q/R is written \mathbf{k}_{\mathrm{phys}}=\mathbf q/a.

\subsection{Newtonian Theory of Small Fluctuations}\label{sec:15.9}

\begin{sloppypar}
We here calculate the behavior of small fluctuations, using the Newtonian
equations \eqref{eq:15.8.1}--\eqref{eq:15.8.4}, but now taking into account
the expansion of the universe. As remarked at the end of Section~\ref{sec:15.1},
we can safely employ Newtonian mechanics to deal with astronomical problems in
which the energy density is dominated by nonrelativistic particles, so that
\(p\ll\rho\), and in which the linear scales involved are small compared with
the characteristic scale of the
universe.\hyperref[ch15-ref-153]{\textsuperscript{153}}
\end{sloppypar}

As shown in Section~\ref{sec:15.1}, there is a simple spatially uniform
solution of Eqs.~\eqref{eq:15.8.1}--\eqref{eq:15.8.4}, with
\begin{align}
  \rho
  &=
  \rho_0\left[\frac{a_0}{a(t)}\right]^3,
  \tag{15.9.1}\label{eq:15.9.1}\\
  \mathbf v
  &=
  \mathbf r\,\frac{\dot a(t)}{a(t)}
  =H\mathbf r,
  \tag{15.9.2}\label{eq:15.9.2}\\
  \mathbf g
  &=
  -\mathbf r\,\frac{4\pi G\rho}{3},
  \tag{15.9.3}\label{eq:15.9.3}
\end{align}
where \(a(t)\) is a scale factor satisfying the differential equation
\begin{equation}
  \dot a^2+k=\frac{8\pi G\rho a^2}{3},
  \tag{15.9.4}\label{eq:15.9.4}
\end{equation}
or equivalently
\begin{equation}
  \frac{\ddot a}{a}=-\frac{4\pi G\rho}{3}.
  \tag{15.9.5}\label{eq:15.9.5}
\end{equation}

We now seek a perturbed solution, by adding to the ``zero-order'' solution
\eqref{eq:15.9.1}--\eqref{eq:15.9.3} the small perturbations
\(\rho_1\), \(\mathbf v_1\), and \(\mathbf g_1\). The hydrodynamic
equations \eqref{eq:15.8.1}--\eqref{eq:15.8.4} then give, to first order
in these perturbations,
\begin{align}
  \dot\rho_1
  +3H\rho_1
  +H(\mathbf r\cdot\boldsymbol{\nabla})\rho_1
  +\rho\,\boldsymbol{\nabla}\cdot\mathbf v_1
  &=0,
  \tag{15.9.6}\label{eq:15.9.6}\\
  \dot{\mathbf v}_1
  +H\mathbf v_1
  +H(\mathbf r\cdot\boldsymbol{\nabla})\mathbf v_1
  &=
  -\frac1\rho\boldsymbol{\nabla}p_1+\mathbf g_1,
  \tag{15.9.7}\label{eq:15.9.7}\\
  \boldsymbol{\nabla}\times\mathbf g_1
  &=0,
  \tag{15.9.8}\label{eq:15.9.8}\\
  \boldsymbol{\nabla}\cdot\mathbf g_1
  &=-4\pi G\rho_1.
  \tag{15.9.9}\label{eq:15.9.9}
\end{align}
Also, since these are for the moment supposed to be adiabatic fluctuations,
the pressure perturbation is given by
\begin{equation}
  p_1=v_s^2\rho_1,
  \tag{15.9.10}\label{eq:15.9.10}
\end{equation}
where \(v_s\) is the speed of sound. (Here and below, \(\rho\) denotes the
mass density \eqref{eq:15.9.1} in the unperturbed solution.)

The equations \eqref{eq:15.9.6}--\eqref{eq:15.9.10} are spatially
homogeneous, so we expect to find plane-wave solutions. Indeed, solutions can
be found with the spatial dependence
\begin{equation}
  \rho_1(\mathbf r,t)
  =
  \rho_1(t)
  \exp\left[
    \frac{\ii\,\mathbf r\cdot\mathbf q}{a(t)}
  \right],
  \tag{15.9.11}\label{eq:15.9.11}
\end{equation}
and likewise for \(\mathbf v_1\) and \(\mathbf g_1\). (The appearance of
the factor \(1/a\) in the exponential means that the wavelength of these
modes is stretched out by the expansion of the universe, as anticipated in
the last section.) Equations~\eqref{eq:15.9.6}--\eqref{eq:15.9.10} now
become coupled ordinary differential equations:
\begin{align}
  \dot\rho_1+3H\rho_1
  +\ii a^{-1}\rho\,\mathbf q\cdot\mathbf v_1
  &=0,
  \tag{15.9.12}\label{eq:15.9.12}\\
  \dot{\mathbf v}_1+H\mathbf v_1
  &=
  -\frac{\ii v_s^2}{\rho a}\,\mathbf q\,\rho_1+\mathbf g_1,
  \tag{15.9.13}\label{eq:15.9.13}\\
  \mathbf q\times\mathbf g_1
  &=0,
  \tag{15.9.14}\label{eq:15.9.14}\\
  \ii\mathbf q\cdot\mathbf g_1
  &=-4\pi G a\rho_1.
  \tag{15.9.15}\label{eq:15.9.15}
\end{align}
The ``field equations'' \eqref{eq:15.9.14} and \eqref{eq:15.9.15} have
the obvious solution
\begin{equation}
  \mathbf g_1
  =
  \frac{4\pi\ii G\rho_1a}{q^2}\,\mathbf q.
  \tag{15.9.16}\label{eq:15.9.16}
\end{equation}
To solve the equations of motion, it is convenient to decompose
\(\mathbf v_1\) into parts perpendicular and parallel to \(\mathbf q\),
\begin{equation}
  \mathbf v_1(t)
  =
  \mathbf v_{1\perp}(t)+\ii\mathbf q\,\epsilon(t),
  \tag{15.9.17}\label{eq:15.9.17}
\end{equation}
where
\[
  \mathbf q\cdot\mathbf v_{1\perp}=0,
  \qquad
  \epsilon=-\frac{\ii\mathbf q\cdot\mathbf v_1}{q^2}.
\]
It is also convenient to express \(\rho_1\) in terms of a fractional change
in density \(\delta\):
\begin{equation}
  \rho_1(t)
  =
  \rho(t)\delta(t)
  =
  \rho_0\left[\frac{a_0}{a(t)}\right]^3\delta(t).
  \tag{15.9.18}\label{eq:15.9.18}
\end{equation}
Then Eq.~\eqref{eq:15.9.13} splits into two uncoupled equations,
\begin{align}
  \dot{\mathbf v}_{1\perp}+H\mathbf v_{1\perp}
  &=0,
  \tag{15.9.19}\label{eq:15.9.19}\\
  \dot\epsilon+H\epsilon
  &=
  -\left(
    \frac{v_s^2}{a}
    -\frac{4\pi G\rho a}{q^2}
  \right)\delta,
  \tag{15.9.20}\label{eq:15.9.20}
\end{align}
and Eq.~\eqref{eq:15.9.12} simplifies to
\begin{equation}
  \dot\delta=\frac{q^2}{a}\epsilon.
  \tag{15.9.21}\label{eq:15.9.21}
\end{equation}

Inspection of Eqs.~\eqref{eq:15.9.19}--\eqref{eq:15.9.21} shows that
there are two quite different types of normal mode here. The \emph{rotational
modes}, described by \(\mathbf v_{1\perp}\), simply decay as \(1/a\):
\begin{equation}
  \mathbf v_{1\perp}(t)\propto a^{-1}(t).
  \tag{15.9.22}\label{eq:15.9.22}
\end{equation}
On the other hand, the \emph{compressional modes}, described by
\(\epsilon\) and \(\delta\), have a more interesting time dependence. Using
Eq.~\eqref{eq:15.9.21} to eliminate \(\epsilon\) in
Eq.~\eqref{eq:15.9.20}, we find
\begin{equation}
  \ddot\delta
  +2H\dot\delta
  +\left(
    \frac{v_s^2q^2}{a^2}-4\pi G\rho
  \right)\delta
  =0.
  \tag{15.9.23}\label{eq:15.9.23}
\end{equation}
Note that this goes over to the Jeans dispersion relation \eqref{eq:15.8.6}
if we set \(a\) constant and define the physical wavevector as
\(\mathbf{k}_{\mathrm{phys}}=\mathbf q/a\).
Equation~\eqref{eq:15.9.23} is the fundamental differential equation that
governs the growth or decay of gravitational condensations in an expanding
universe.

The above Newtonian theory becomes applicable at the onset of the
matter-dominated era, when the energy density of radiation drops below the
rest-mass density, so that \(p\ll\rho\). Unfortunately,
Eq.~\eqref{eq:15.9.23} is a little too complicated to allow a solution in
closed form that would be valid throughout the whole of the matter-dominated
era. We can, however, answer the most interesting questions about the behavior
of \(\delta(t)\) by solving Eq.~\eqref{eq:15.9.23} in a number of special
cases.

\subsubsection*{(A) Zero Pressure Solutions}

According to the general picture outlined in the last section, a galaxy is
supposed to grow out of the small fluctuations present at the time of hydrogen
recombination, which are left over from the previous phase of damped acoustic
oscillation. The most important problem facing us is how much such a
protogalactic fluctuation could have grown from the time of recombination until
the present. Or, to put this in terms relevant to observations, how large would
a fluctuation have to be at the time of recombination to have a chance of
growing into a galaxy by the present time?

To answer these questions we can simplify Eq.~\eqref{eq:15.9.23} by
neglecting the pressure term \(v_s^2q^2/a^2\). This term will be negligible
compared with the gravitational term \(4\pi G\rho\) if the physical wave
number \(\lvert\mathbf{k}_{\mathrm{phys}}\rvert=\lvert\mathbf q\rvert/a\)
is much less than the Jeans wave number \eqref{eq:15.8.8} or, equivalently,
if the fluctuation mass \eqref{eq:15.8.12} is much greater than the Jeans
mass \eqref{eq:15.8.13}. We saw in the last section that the Jeans mass is
of order \(10^6M_\odot\) to \(3\times10^6M_\odot\) immediately after
recombination, and drops as \(a^{-3/2}\) thereafter, so a galactic mass of
order \(10^{11}M_\odot\) will certainly be much greater than the Jeans mass
once the hydrogen recombination is over.

In order to carry the solution of Eq.~\eqref{eq:15.9.23} forward to the
present time, it will be necessary to use the parametric formulas for \(a(t)\)
and \(\rho(t)\) derived in Section~\ref{sec:15.3}. For positive, zero, and
negative curvature, these are:
\[
\begin{array}{c@{\qquad}l}
k=+1&
\begin{aligned}
  \frac{a(t)}{a_0}
  &=q_0(2q_0-1)^{-1}(1-\cos\theta),\\
  H_0t
  &=q_0(2q_0-1)^{-3/2}(\theta-\sin\theta),\\
  \rho
  &=
  \frac{3H_0^2(2q_0-1)^3}
       {4\pi Gq_0^2(1-\cos\theta)^3};
\end{aligned}\\[2.8em]
k=0&
\begin{aligned}
  \frac{a(t)}{a_0}
  &=\left(\frac{3H_0t}{2}\right)^{2/3},\\
  \rho&=(6\pi Gt^2)^{-1};
\end{aligned}\\[2.2em]
k=-1&
\begin{aligned}
  \frac{a(t)}{a_0}
  &=q_0(1-2q_0)^{-1}(\cosh\Psi-1),\\
  H_0t
  &=q_0(1-2q_0)^{-3/2}(\sinh\Psi-\Psi),\\
  \rho
  &=
  \frac{3H_0^2(1-2q_0)^3}
       {4\pi Gq_0^2(\cosh\Psi-1)^3}.
\end{aligned}
\end{array}
\]
With the pressure term neglected, the differential equation
\eqref{eq:15.9.23} now takes the following forms:
\begin{align}
  (1-\cos\theta)\frac{\dd^2\delta}{\dd\theta^2}
  +\sin\theta\frac{\dd\delta}{\dd\theta}
  -3\delta
  &=0
  &&(k=+1),
  \tag{15.9.24}\label{eq:15.9.24}\\
  \ddot\delta+\frac4{3t}\dot\delta-\frac2{3t^2}\delta
  &=0
  &&(k=0),
  \tag{15.9.25}\label{eq:15.9.25}\\
  (\cosh\Psi-1)\frac{\dd^2\delta}{\dd\Psi^2}
  +\sinh\Psi\frac{\dd\delta}{\dd\Psi}
  -3\delta
  &=0
  &&(k=-1).
  \tag{15.9.26}\label{eq:15.9.26}
\end{align}
In each case there are two independent solutions, which we shall call
\(\delta_+\) and \(\delta_-\):
\begin{align}
  \delta_+
  &\propto
  -\frac{3\theta\sin\theta}{(1-\cos\theta)^2}
  +\frac{5+\cos\theta}{1-\cos\theta},
  &&(k=+1),
  \tag{15.9.27}\label{eq:15.9.27}\\
  \delta_-
  &\propto
  \frac{\sin\theta}{(1-\cos\theta)^2},
  &&(k=+1),
  \tag{15.9.28}\label{eq:15.9.28}\\
  \delta_+
  &\propto t^{2/3},
  &&(k=0),
  \tag{15.9.29}\label{eq:15.9.29}\\
  \delta_-
  &\propto t^{-1},
  &&(k=0),
  \tag{15.9.30}\label{eq:15.9.30}\\
  \delta_+
  &\propto
  -\frac{3\Psi\sinh\Psi}{(\cosh\Psi-1)^2}
  +\frac{5+\cosh\Psi}{\cosh\Psi-1},
  &&(k=-1),
  \tag{15.9.31}\label{eq:15.9.31}\\
  \delta_-
  &\propto
  \frac{\sinh\Psi}{(\cosh\Psi-1)^2},
  &&(k=-1).
  \tag{15.9.32}\label{eq:15.9.32}
\end{align}
In all three cases, the solutions \(\delta_+\) and \(\delta_-\) go over to
\(t^{2/3}\) and \(t^{-1}\), respectively, for \(a(t)\ll a_0\), so that if
a disturbance starts with comparable amplitudes for the \(\delta_+\) and
\(\delta_-\) modes at recombination, then it will soon be almost purely in
the \(\delta_+\) mode. We therefore concentrate entirely on the
\(\delta_+\) mode from now on.

The disturbance is supposed here to start at a time \(t_R\) corresponding to
the large red shift
\[
  1+z_R
  \equiv
  \frac{a(t_0)}{a(t_R)}
  \approx
  \frac{4000\,\mathrm K}{2.7\,\mathrm K}
  \approx1500.
\]
The initial values of the parameters \(\theta\), \(t\), or \(\Psi\) for
\(k=+1\), \(k=0\), or \(k=-1\) are then
\[
  \theta_R
  \approx
  \left[
    \frac{2(1-\cos\theta_0)}{1+z_R}
  \right]^{1/2},
  \qquad
  t_R\approx t_0(1+z_R)^{-3/2},
  \qquad
  \Psi_R
  \approx
  \left[
    \frac{2(\cosh\Psi_0-1)}{1+z_R}
  \right]^{1/2}.
\]
Hence the density contrast \(\delta\) could at most grow by an amplification
factor
\begin{equation}
  A_0
  \equiv
  \frac{\delta_+(t_0)}{\delta_+(t_R)},
  \tag{15.9.33}\label{eq:15.9.33}
\end{equation}
given by
\begin{equation}
  A_0
  =
  \begin{cases}
    \displaystyle
    \frac{5(1+z_R)}{(1-\cos\theta_0)^3}
    \left[
      -3\theta_0\sin\theta_0
      +(1-\cos\theta_0)(5+\cos\theta_0)
    \right],
    &k=+1,\\[1.2em]
    1+z_R,
    &k=0,\\[0.8em]
    \displaystyle
    \frac{5(1+z_R)}{(\cosh\Psi_0-1)^3}
    \left[
      -3\Psi_0\sinh\Psi_0
      +(\cosh\Psi_0-1)(\cosh\Psi_0+5)
    \right],
    &k=-1.
  \end{cases}
  \tag{15.9.34}\label{eq:15.9.34}
\end{equation}
The parameters \(\theta_0\) and \(\Psi_0\) can be expressed in terms of
\(q_0\) by using Eq.~\eqref{eq:15.3.10} or \eqref{eq:15.3.19}. We then
find that \(A_0\) is a monotonically increasing function of \(q_0\) for
\(q_0>0\), rising from \(A_0\approx5q_0(1+z_R)\) for \(q_0\ll\tfrac12\),
to \(A_0=1+z_R\) at \(q_0=\tfrac12\), to
\(A_0=1.45(1+z_R)\) at \(q_0=1\), and approaching the upper bound
\(A_0=5(1+z_R)\) for \(q_0\gg1\). With \(q_0\) between 0.014 and 2 and
\(1+z_R=1500\), a small fluctuation would have grown from the time of
hydrogen recombination until the present by a factor \(A_0\) between 100
and 3000.

The condensations observed in the present universe cannot be considered
``small disturbances.'' For instance, the mass density in a typical cluster
of galaxies is of order \(10^{-28}\,\mathrm{g/cm^3}\), an order of magnitude
greater than the maximum likely value for the universe as a whole, and the mass
density within a galaxy is of course even greater. The simple linear
instability theory described above is therefore not applicable to the whole
history of inhomogeneities up to the present moment. However, it seems
reasonable to suppose that the present strong condensations grew out of small
disturbances, so that a necessary (if perhaps not sufficient) condition for
their formation is that the perturbation \(\delta_+(t)\) calculated in linear
stability theory should have become of order unity at some time before the
present. This then sets a lower limit on the magnitude of the initial
disturbance at the time of recombination:
\begin{equation}
  \lvert\delta_+(t_R)\rvert\gtrsim\frac1{A_0},
  \tag{15.9.35}\label{eq:15.9.35}
\end{equation}
so that \(\lvert\delta_+(t_R)\rvert\) should be greater than about
\(10^{-2}\) to \(3\times10^{-4}\), depending on the value of \(q_0\). In
order to say how much greater, it would be necessary to know the time at which
the disturbance reached the beginning of the nonlinear regime, with
\(\lvert\delta_+(t)\rvert\approx1\). According to
Weymann,\hyperref[ch15-ref-154]{\textsuperscript{154}} the observed binding
energies of galaxies indicates that this must have occurred after a time of
order \(7\times10^7\) years; if \(\delta_+\) reached unity this early, then
it must have been already quite large at the time of recombination. On the
other hand, if the concentration of quasi-stellar object red shifts at
\(z\approx2\) marks the onset of galaxy formation, then
Eq.~\eqref{eq:15.9.35} provides a reasonably good estimate of the actual
value of \(\lvert\delta_+(t_R)\rvert\). As remarked at the end of the last
section, the protogalactic fluctuations at the time of recombination would
produce fractional fluctuations in the observed microwave background
temperature equal to \(\lvert\delta_+(t_R)\rvert/3\) over angles of order
\(30''\), provided that Thomson scattering during or after recombination does
not smooth out the
fluctuations.\hyperref[ch15-ref-152a]{\textsuperscript{152a}} Even if the
nonlinear regime was reached rather recently, \(\Delta T_\gamma/T_\gamma\)
would be of order \(3\times10^{-3}\) to \(10^{-4}\), which should be
observable.

\subsubsection*{(B) Zero Curvature Solutions}

It is also interesting to consider the behavior of the solutions when the
pressure term \(v_s^2q^2/a^2\) in Eq.~\eqref{eq:15.9.23} is not neglected,
particularly in order to locate the precise dividing line between stable and
unstable fluctuations. In order to obtain exact solutions, it is necessary now
to restrict our attention to early times, when \(a(t)\ll a_0\), so that the
terms \(\dot a^2\) and \(8\pi\rho Ga^2/3\) in Eq.~\eqref{eq:15.9.4} are
much larger than unity, and even for \(k=\pm1\) we can use the \(k=0\)
formulas for \(a\) and \(\rho\):
\begin{align}
  a&\propto t^{2/3},
  \tag{15.9.36}\label{eq:15.9.36}\\
  \rho&=(6\pi Gt^2)^{-1}.
  \tag{15.9.37}\label{eq:15.9.37}
\end{align}
(This is not much of a limitation, because for recent times, when these
formulas may not be valid, the Jeans mass is so small that its precise value is
of little interest.) For a general specific heat ratio \(\gamma\), the
pressure varies as \(\rho^\gamma\), and the speed of sound is
\begin{equation}
  v_s
  =
  \left(\frac{\gamma p}{\rho}\right)^{1/2}
  \propto\rho^{(\gamma-1)/2}
  \propto t^{1-\gamma}.
  \tag{15.9.38}\label{eq:15.9.38}
\end{equation}
Hence Eq.~\eqref{eq:15.9.23} reads here
\begin{equation}
  \ddot\delta
  +\frac4{3t}\dot\delta
  +\left(
    \frac{\Lambda^2}{t^{2\gamma-2/3}}
    -\frac2{3t^2}
  \right)\delta
  =0,
  \tag{15.9.39}\label{eq:15.9.39}
\end{equation}
where \(\Lambda^2\) is the constant
\begin{equation}
  \Lambda^2
  \equiv
  \frac{t^{2\gamma-2/3}v_s^2q^2}{a^2}.
  \tag{15.9.40}\label{eq:15.9.40}
\end{equation}
The solutions of Eq.~\eqref{eq:15.9.39} for \(\gamma>4/3\) are
\begin{equation}
  \delta_\pm
  \propto
  t^{-1/6}
  J_{\mp5/(6\nu)}
  \left(\frac{\Lambda t^{-\nu}}{\nu}\right),
  \tag{15.9.41}\label{eq:15.9.41}
\end{equation}
where \(J\) is the usual Bessel function, and
\begin{equation}
  \nu\equiv\gamma-\frac43>0.
  \tag{15.9.42}\label{eq:15.9.42}
\end{equation}
The Bessel functions oscillate for \(t\ll\Lambda^{1/\nu}\), whereas for
\(t\gg\Lambda^{1/\nu}\) the solutions behave like
\begin{equation}
  \delta_\pm\propto t^{-1/6\pm5/6}.
  \tag{15.9.43}\label{eq:15.9.43}
\end{equation}
By using Eqs.~\eqref{eq:15.9.37} and \eqref{eq:15.9.40}, the condition
\(t>\Lambda^{1/\nu}\) for growth in the \(\delta_+\) mode can be written
% SOURCE ERRATUM: the scan prints greater-than signs here, but (15.9.23) and
% the later criterion after (15.9.46) require pressure support below the
% Jeans threshold.
\[
  \frac{v_s^2q^2}{a^2}\lesssim6\pi G\rho,
\]
which is substantially the same as the Jeans condition
\(v_s^2k_{\mathrm{phys}}^2<4\pi G\rho\).

The solutions \eqref{eq:15.9.41} will apply after recombination, with
\(\gamma\approx5/3\). Also, for a relatively high present density, there is
an appreciable period before recombination when the total energy density is
dominated by hydrogen rest mass, but the pressure is dominated by radiation.
The cosmic medium then behaves like a nonrelativistic fluid with
\(\gamma=4/3\), so that Eq.~\eqref{eq:15.9.39} reads
\begin{equation}
  \ddot\delta
  +\frac4{3t}\dot\delta
  +\left(\Lambda^2-\frac23\right)\frac{\delta}{t^2}
  =0,
  \tag{15.9.44}\label{eq:15.9.44}
\end{equation}
with
\begin{equation}
  \Lambda^2
  \equiv
  \frac{t^2v_s^2q^2}{a^2}
  =
  \frac{v_s^2q^2}{6\pi G\rho a^2}.
  \tag{15.9.45}\label{eq:15.9.45}
\end{equation}
The solutions of Eq.~\eqref{eq:15.9.44} are quite simple:
\begin{equation}
  \delta_\pm\propto t^\alpha,
  \qquad
  \alpha=-\frac16\pm\left(\frac{25}{36}-\Lambda^2\right)^{1/2}.
  \tag{15.9.46}\label{eq:15.9.46}
\end{equation}
Both solutions undergo a gently damped oscillation for
\(\Lambda>5/6\), and decay for \(5/6>\Lambda>\sqrt{2/3}\), whereas for
\(\Lambda<\sqrt{2/3}\), \(\delta_+\) grows and \(\delta_-\) decays. The
condition for growth in the \(\delta_+\) mode is then
\[
  \frac{v_s^2q^2}{6\pi G\rho a^2}<\frac23,
\]
which is precisely the same as Jeans's condition
\(v_s^2k_{\mathrm{phys}}^2<4\pi G\rho\).

It is not difficult to incorporate dissipative effects in this formalism. The
most interesting case here is the damping owing to a finite photon mean free
path during the matter-dominated part of ``Phase B,'' that is, during the
period prior to recombination when the radiation density \(uT^4\) is much
less than the matter density \(nm_{\mathrm H}\) and the Jeans mass is much
greater than a galactic mass. According to
Eqs.~\eqref{eq:15.8.27}--\eqref{eq:15.8.30}, the effects of viscosity are
negligible here compared with the effects of heat conduction, so dissipative
effects can be taken into
account\hyperref[ch15-ref-155]{\textsuperscript{155}} by using the equation
of state to express the pressure perturbation \(p_1\) in
Eq.~\eqref{eq:15.9.7} in terms of the temperature and mass-density
perturbations \(T_1\) and \(\rho_1\), and supplementing
Eqs.~\eqref{eq:15.9.6}--\eqref{eq:15.9.9} with the usual nonrelativistic
equation of radiative heat transfer. We need not go into this further here,
since heat conduction as well as viscosity will be incorporated into the
general-relativistic theory described in the next section.
